
\begin{table}[H]
   \caption{\label{tab:pretrend_robustness} \textbf{\emph{Parto Adequado}: covariate and treated-trend robustness}}
   \centering
\small\setlength{\tabcolsep}{5pt}\resizebox{\ifdim\width>\linewidth \linewidth\else\width\fi}{!}{%
   \begin{tabular}{lcccc}
      \toprule
      Dependent Variable: & \multicolumn{4}{c}{Cesarean rate}\\
                                   & Baseline        & + Covariates     & + Treated trend      & + Both \\
      Model:                       & (1)             & (2)              & (3)                  & (4)\\
      \midrule
      \emph{Variables}\\
      Treated $\times$ Post (2017+) & -0.0659$^{***}$ & -0.0508$^{***}$  & 0.0029               & 0.0042\\
                                   & (0.0146)        & (0.0145)         & (0.0099)             & (0.0099)\\   
\addlinespace[2pt]
      \midrule
      Municipality fixed effects & Yes             & Yes              & Yes                  & Yes\\  
      Year fixed effects & Yes             & Yes              & Yes                  & Yes\\  
      \midrule
      \emph{Fit statistics}\\
      Observations                 & 7,660           & 7,218            & 7,660                & 7,218\\  
\bottomrule
   \end{tabular}
}
   
   \par \raggedright 
   \footnotesize\textit{Notes:} SINASC for-profit cesarean rate, municipality-year 2010--2024, weighted by births. Column 1 is the baseline difference-in-differences; column 2 adds time-varying municipal covariates; column 3 adds a treated-cohort linear time trend; column 4 adds both. Covariates leave the differential pre-trend intact (joint 2010--2015 test still rejects, $p<0.01$); the treated linear trend absorbs the 2017 ``break'' entirely, confirming it is a pre-existing trend, not a treatment effect. Standard errors, clustered by municipality, are reported in parentheses. \footnotesize Significance levels: *** p$<$0.01, ** p$<$0.05, * p$<$0.10.
\end{table}


